\section{Step-AUC Analysis}
\label{sec:results_step_auc}

Using the sample-efficiency metric defined in Section~\ref{subsec:methodology_sample_efficiency_metrics}, step-AUC measures average evaluation return over the full training budget.
For a method with mean evaluation return $\bar{G}(s_i)$ at checkpoint step $s_i$, the normalized step-AUC is computed as
\begin{equation}
    \mathrm{AUC}_{\mathrm{step}}
    = \frac{1}{s_K-s_0}\sum_{i=1}^{K}
      \frac{\bar{G}(s_i)+\bar{G}(s_{i-1})}{2}(s_i-s_{i-1}),
    \label{eq:results_step_auc}
\end{equation}
where $s_0=0$ and $s_K$ is the final evaluation step.
This metric rewards methods that learn early and maintain high return, rather than methods that improve only at the final checkpoint.

Table~\ref{tab:results_step_auc_glpe_vs_baselines} reports step-AUC for GLPE, GLPE without gating, and the strongest baseline per environment.
The best baseline is selected independently for each environment from PPO, ICM, RND, RIDE, and RIAC.

\begin{table}[!htbp]
\centering
\small
\begin{fitwidthtabular}{lccc}
\hline
Environment & GLPE & GLPE (no gate) & Best baseline \\
\hline
\textsc{MountainCar-v0} & $-107.0\;[-114.8,-101.8]$ & $-111.3\;[-128.7,-101.4]$ & RIDE: $-113.7\;[-134.2,-101.9]$ \\
\textsc{BipedalWalker-v3} & $249.0\;[223.9,270.7]$ & $255.2\;[230.0,275.6]$ & ICM: $257.3\;[229.4,278.3]$ \\
\textsc{Ant-v5} & $3{,}402\;[3{,}067,3{,}708]$ & $3{,}402\;[3{,}067,3{,}707]$ & RND: $3{,}565\;[3{,}202,3{,}876]$ \\
\textsc{HalfCheetah-v5} & $5{,}005\;[4{,}708,5{,}316]$ & $4{,}998\;[4{,}700,5{,}296]$ & RIDE: $5{,}208\;[5{,}019,5{,}410]$ \\
\textsc{Humanoid-v5} & $1{,}583\;[551,3{,}360]$ & $1{,}665\;[769,2{,}921]$ & ICM: $1{,}781\;[829,2{,}797]$ \\
\hline
\end{fitwidthtabular}
\caption[Step-AUC by environment and method]{Step-AUC of deterministic evaluation return versus environment steps. Values are normalized by the corresponding environment step budget. Brackets show approximate 95\% bootstrap uncertainty intervals obtained from the checkpoint confidence bands. Higher values are better.}
\label{tab:results_step_auc_glpe_vs_baselines}
\end{table}

GLPE has the highest step-AUC on \textsc{MountainCar-v0}, which is consistent with its rapid and reliable discovery of the sparse-reward solution.
On \textsc{BipedalWalker-v3}, GLPE without gating is close to ICM, differing by about $2.1$ return points in mean step-AUC.
On \textsc{Ant-v5} and \textsc{HalfCheetah-v5}, the strongest baselines have a moderate advantage, but the GLPE variants remain within the same broad performance range.
On \textsc{Humanoid-v5}, ICM has the highest mean step-AUC, but the confidence intervals for GLPE, GLPE without gating, and ICM overlap substantially.

The step-AUC results indicate that GLPE is not merely a final-checkpoint effect in the sparse-reward task.
Its advantage on \textsc{MountainCar-v0} is accumulated throughout training.
In contrast, the smaller gaps in the locomotion environments show that the proposed intrinsic reward improves neither all training phases nor all task families uniformly.
This supports evaluating GLPE as a targeted exploration method rather than as a replacement for all intrinsic-reward approaches.